\chapter*{\large ВВЕДЕНИЕ}  
\addcontentsline{toc}{chapter}{ВВЕДЕНИЕ}
Люди на протяжении своего существования придумывали различные истории, преследуюя множество целей: попытка объяснить происходящее, для развлечения, обучения и др. Со временем оформились отдельные жанры литературы, которые основаны на сочинении различных миров - фантастика. Логично, что для большего "погружения" в выдуманный мир требуется нечто большее, чем просто рассказать о нем. Кто-то придумывает, на каких выдуманных языках разговаривают в их мирах, в каких богов. Но чаще всего это различные илюстрации -  изображения существ, карты мира. И создание этих карт  представляет сложную задачу: Определить, сколько и каких материков, острово должно быть? какой формы они должны быть?

Спустя время появились Вычислительные машины, которые позволили упростить данный процесс. Но куда больше облегчили этот процесс алгоритмы: Шум Перлина(1982,Перлин), Симплексный шум(2001, Перлин), Diamond Square алгоритм( Фурнье, Фуссель и Карпентер, 1982), использование клеточных автоматов.